% -*- compile-command: "pdflatex -shell-escape Hendrix-comprog-reference.tex" -*-

\documentclass[10pt]{book}

\usepackage[margin=1in]{geometry}
\usepackage{fancyhdr}

\pagestyle{fancy}

% \lhead{\leftmark}
% \rhead{\rightmark}

% \lfoot{Hendrix Programming Team Reference}
% \rfoot{\thepage}

% TODO: make fancy header with page number and section etc.

\usepackage[T1]{fontenc}
\usepackage[utf8]{inputenc}
\usepackage{hyperref}
\usepackage{xcolor}

\hypersetup{
    colorlinks,
    linkcolor={green!50!black},
    citecolor={blue!50!black},
    urlcolor={blue!80!black}
}

\usepackage{titlepic}
\usepackage{graphicx}
\graphicspath{{images/}}

\usepackage[newfloat]{minted}
\usepackage{etoolbox}
\usepackage{mdframed}
\usepackage{booktabs}
\usepackage{xspace}
\usepackage{pifont}

%%%%%%%%%%%%%%%%%%%%%%%%%%%%%%%%%%%%%%%%%%%%%%%%%%
%% Inline code

\newcommand{\code}[1]{\inputminted[fontsize=\normalsize]{java}{code/#1}}
\newcommand{\pycode}[1]{\inputminted[fontsize=\normalsize]{python}{code/#1}}

%%%%%%%%%%%%%%%%%%%%%%%%%%%%%%%%%%%%%%%%%%%%%%%%%%
%% Links

\newcommand*{\fulllink}[1]{\hyperref[{#1}]{\nameref*{#1}~(\S\ref*{#1}, page~\pageref*{#1})}}
\newcommand*{\link}[1]{\hyperref[{#1}]{(\S\ref*{#1}, page~\pageref*{#1})}}

%%%%%%%%%%%%%%%%%%%%%%%%%%%%%%%%%%%%%%%%%%%%%%%%%%
%% Kattis

\newcommand{\kattislogo}{\raisebox{-0.2em}{\includegraphics[height=0.9\baselineskip]{Kattis}}}

\newcommand{\kattis}[1]{%
  \begin{mdframed}
  \kattislogo
  \def\nextitem{\def\nextitem{, }}
  \renewcommand*{\do}[1]{\nextitem\kattislink{##1}}
  \docsvlist{#1}
  \end{mdframed}
}

\newcommand{\kattislink}[1]{\href{https://open.kattis.com/problems/#1}{\texttt{#1}}}

\newcommand{\inlinekattis}[1]{\kattislogo\ \kattislink{#1}}

%%%%%%%%%%%%%%%%%%%%%%%%%%%%%%%%%%%%%%%%%%%%%%%%%%
%% Todos

\newif\iftodos
\todostrue
\newcommand{\todo}[1]{\iftodos\textcolor{red}{[TODO: #1]}\fi}

%%%%%%%%%%%%%%%%%%%%%%%%%%%%%%%%%%%%%%%%%%%%%%%%%%
%% Warning

\newenvironment{warning}
{\par\begin{mdframed}[linewidth=2pt,linecolor=red]%
    \begin{list}{}{\leftmargin=1cm
        \labelwidth=\leftmargin}\item[\Large\ding{43}]}
    {\end{list}\end{mdframed}\par}

%%%%%%%%%%%%%%%%%%%%%%%%%%%%%%%%%%%%%%%%%%%%%%%%%%
%% Misc

\newcommand{\eg}{\emph{e.g.}\xspace}
\newcommand{\ie}{\emph{i.e.}\xspace}

%%%%%%%%%%%%%%%%%%%%%%%%%%%%%%%%%%%%%%%%%%%%%%%%%%%%%%%%%%%%%%%%%%%%%%

\begin{document}

\title{Hendrix Programming Team Reference}
\titlepic{\includegraphics[width=3in]{Hendrix-logo}}
\maketitle

\tableofcontents
\newpage

\chapter{Limits}

As a rule of thumb, you should assume about $10^8$ (= 100 million)
operations per second.  If you can think of a straightforward brute
force solution to a problem, you should check whether it is likely to
fit within the time limit; if so, go for it!  Some problems are
explicitly written to see if you will recognize this.  If a brute
force solution won't fit, the input size can help guide you to search
for the right algorithm running time.

Example: suppose a problem requires you to find the length of a
shortest path in a weighted graph.
\begin{itemize}
\item If the graph has $|V| = 400$ vertices, you should use
  Floyd-Warshall \link{sec:floydwarshall}: it is the easiest to code and takes $O(V^3)$ time
  which should be good enough.
\item If the graph has $|V| = 4000$ vertices, especially if it doesn't
  have all possible edges, you can use Dijkstra's algorithm
  \link{sec:dijkstra}, which is $O(E \log V)$.
\item If the graph has $|V| = 10^5$ vertices, you should look for some
  special property of the graph which allows you to solve the problem
  in $O(V)$ or $O(V \log V)$ time---for example, perhaps the graph is
  a tree \link{sec:graph-basics}, so you can run a DFS \link{sec:dfs}
  to find a unique path and then add up the weights.  An input size of
  $10^5$ is a common sign that you are expected to use an $O(n \lg n)$
  or $O(n)$ algorithm---it's big enough to make $O(n^2)$ too slow but
  not so big that the time to do I/O makes a big difference.
\end{itemize}

\begin{table*}[h]\centering
  \renewcommand{\arraystretch}{1.3}
  \begin{tabular}{@{}lll@{}}\toprule
    $n$ & Worst viable running time & Example \\
    \midrule
    $11$ & $O(n!)$ & Generating all permutations \link{sec:combinatorics} \\
    $25$ & $O(2^n)$ & Generating all subsets \link{chap:bittricks} \\
    $100$ & $O(n^4)$ & Some brute force algorithms \\
    $400$ & $O(n^3)$ & Floyd-Warshall \link{sec:floydwarshall} \\
    $10^4$ & $O(n^2)$ & Testing all pairs \\
    $10^6$ & $O(n \lg n)$ & BFS/DFS; sort+greedy \\
    \bottomrule
  \end{tabular}
\end{table*}

\kattis{bing, transportationplanning, dancerecital, prozor,
  rectanglesurrounding, weakvertices}

\begin{itemize}
\item $2^{10} = 1024 \approx 10^3$.
\item One \texttt{int} is 32 bits = 4 bytes. So \emph{e.g.} an array
  of $10^6$ \texttt{int}s requires $< 4$ MB---no big deal since the
  typical memory limit is $1$ GB.  Don't be afraid to make arrays with
  millions of elements!
\item \texttt{int} holds 32 bits; the largest \texttt{int} value is
  \verb|Integer.MAX_VALUE| $ = 2^{31} - 1$, a bit more than
  $2 \cdot 10^9$.
\item \texttt{long} holds 64 bits; the largest
  \texttt{long} value is \verb|Long.MAX_VALUE| $ = 2^{63} - 1$, a bit
  more than $9 \cdot 10^{18}$.
\item If you need larger values, use \fulllink{sec:bigint} or just use
  \fulllink{chap:python}; see \fulllink{sec:combinatorics}.
\end{itemize}

\chapter{Java Reference}

\section{Template}

\code{java/Template.java}

\section{Scanner}

\texttt{Scanner} is relatively slow but should usually be sufficient
for most purposes.  If the input or output is relatively large (> 1MB)
and you suspect the time taken to read or write it may be a hindrance,
you can use \fulllink{sec:fastio}.

\code{java/ScannerExample.java}

\section{String}

The \texttt{String} type can be used in Java to represent sequences of
characters.

\todo{Useful \texttt{String} methods: concatenation, substring, charAt, \dots?}
\todo{Converting between \texttt{String} and \texttt{char[]},
  advantages and disadvantages of each}

\kattis{battlesimulation, bing, connectthedots, itsasecret, shiritori,
suffixarrayreconstruction}

\section{Arrays}

\todo{Basic array template/examples.}
\todo{Useful \texttt{Arrays} methods: sort, copyOf, copyOfRange, fill, binarySearch}

\kattis{falcondive,freefood,traveltheskies}

\section{ArrayList}

\todo{Basic template and examples.  Useful \texttt{ArrayList} methods:
  \texttt{add}, \texttt{get}, \texttt{set}.  Advantages/disadvantages
  compared to arrays.}

\section{Stack}

\todo{Example using \texttt{Stack} class.}
\todo{Mention balanced parentheses, DFS}

\kattis{backspace, islands, pairingsocks, reservoir, restaurant,
  throwns, zagrade}

\section{Queue}

\todo{\texttt{Queue} interface, \texttt{ArrayDeque} class}

\kattis{brexit, coconut, ferryloading4, shuffling}

\section{PriorityQueue}

\todo{Examples of using \texttt{PriorityQueue}.  Show how to
  construct with custom \texttt{Comparator}, eg. using lambda notation}
\todo{Note lack of decreaseKey operation (Dijkstra), use remove + add,
  not as fast}

\kattis{bank, guessthedatastructure, knigsoftheforest}

\section{Set}

\todo{HashSet, TreeSet}

\section{Map}

\todo{HashMap, TreeMap}
\todo{Iterating over keys, values, both (MapEntry)}

\section{BigInteger} \label{sec:bigint}

\todo{Examples.  Useful methods, constructors (gcd, mod, base conversion!).}

\kattis{basicremains}

\section{Sorting}

\todo{Basic template for implementing Comparable}
\todo{Arrays.sort, Collections.sort}
\todo{Sorting with a custom \texttt{Comparator}}
\todo{Include code for basic sorting implementations (in case it's
  useful to code them up explicitly so they can be enhanced with extra
  info): insertion sort, mergesort, quicksort)}

\section{BitSet}

\todo{Basic examples of BitSet use.}

\kattis{primesieve}

\section{Fast I/O} \label{sec:fastio}

Typically ACM ICPC problems are designed so \texttt{Scanner} and
\texttt{System.out.println} are fast enough to read and write the
required input and output within the time limits.  However, these are
relatively slow since they are unbuffered (every single read and write
happens immediately).  Occasionally it can be useful to have faster
I/O; indeed, some problems on Kattis cannot be solved in Java without
using this. \todo{Link to some examples.}

\begin{warning}
  Be sure to call \texttt{flush()} at the end of your program
  or else some output might be lost!
\end{warning}

\code{java/Kattio.java}

\todo{Add \texttt{getLine()} method}

\chapter{Data Structures}

\section{Union-find}

A union-find structure can be used to keep track of a collection of
disjoint sets, with the ability to quickly test whether two items are
in the same set, and to quickly union two given sets into one.  It is
used in Kruskal's Minimum Spanning Tree algorithm \link{sec:kruskal},
and can also be useful on its own.  \texttt{find} and \texttt{union}
both take essentially constant amortized time.

\kattis{drivingrange, islandhopping, kastenlauf, lostmap, minspantree,
  numbersetseasy, treehouses, unionfind, virtualfriends,
  wheresmyinternet}

\code{data-structures/UnionFind.java}

\section{Heaps}

\section{Tries}

\kattis{boggle, heritage, herkabe, phonelist}

\section{Red-black trees}

\section{Segment trees and Fenwick trees}

See \fulllink{sec:range-queries}.

\chapter{Search}

\section{Complete search}

\todo{Complete search aka brute force}

\section{Binary and ternary search}

\todo{Binary search on an array; binary search on unbounded function
  on the integers; binary search on real interval; ternary search}
\todo{Point out \texttt{Arrays.binarySearch}}

\chapter{Graphs}

\section{Graph basics} \label{sec:graph-basics}

\todo{Directed, undirected, weighted, unweighted, self loops, multiple edges}
\todo{characterization of trees}
\todo{New virtual source/sink node trick}

\section{Graph representation}

\todo{Adjacency matrix, adjacency maps.  Edge objects. Implicit graphs.}

\section{BFS}

\todo{Code for BFS with level labelling, parent map.}

\kattis{brexit}

\section{DFS, SCCs, topological sorting} \label{sec:dfs}

\todo{Code for DFS, start/finish labelling, top sorting, Tarjan's SCC algorithm}

\section{Single-source shortest paths (Dijkstra)} \label{sec:dijkstra}

\section{All-pairs shortest paths (Floyd-Warshall)} \label{sec:floydwarshall}

\section{Min spanning trees (Kruskal)} \label{sec:kruskal}

\section{Max flow}

A \emph{flow network} is a directed, weighted graph where the edge
weights (typically integers) are thought of as representing
\emph{capacities} (\eg imagine pipes of varying sizes).  The \emph{max
  flow problem} is to determine, given a flow network, the maximum
possible amount of \emph{flow} which can move through the network
between given source and sink vertices, subject to the constraints
that the flow on any edge is no greater than the capacity, and the sum
of incoming flows equals outgoing flows at every vertex other than the
source or sink.  Flow networks can be used to model a wide variety of
problems.

\todo{Enumerate a few problem types: item assignment; max bipartite
  matching; min cut}

\todo{choose directed/undirected edges carefully!}

\todo{Requires vertices $0 \dots n-1$: either carefully keep track of
  which numbers are for which vertices, or use lookup tables}

\kattis{copsandrobbers,escapeplan,gopher2,guardianofdecency,marblestree,maxflow,mincut,paintball,waif}

Dinitz' Algorithm is probably the best all-around algorithm to use for
solving max flow problems in competitive programming.  It takes
$O(V^2 E)$ in theory (although is often much faster in practice).  In
the special case where we are modelling a bipartite matching problem,
Dinitz' Algorithm reduces to the Hopcroft-Karp algorithm which runs in
$O(E \sqrt{V})$.

\code{flow/Dinitz.java}

\todo{Include a sample solution using a flow network}

\chapter{Dynamic Programming}

\chapter{Strings}

\section{Suffix arrays}

\chapter{Divide \& Conquer}

\section{Counting inversions}

\kattis{excursion, froshweek}

\chapter{Mathematics}

\section{GCD/Euclidean Algorithm}

\section{Fractions}

\section{Primes and factorization}

\todo{Basic primality testing and factorization with trial division.
  Sieving (primes, factors, Euler totient).}

\section{Factorial} \label{sec:factorial}

\kattis{eulersnumber, factstone, howmanydigits, lastfactorialdigit,
  inversefactorial, loworderzeros}

\todo{Computing factorials; size using logs, etc}

\section{Combinatorics} \label{sec:combinatorics}

\kattis{anagramcounting, secretsanta}

\todo{Basic principles of combinatorics.  Code for computing binomial
  coefficients. Multinomial coefficients.}

\todo{mod $10^9 + 7$.}

\begin{warning}
  Remember to use \texttt{long} if you need an answer
  $\bmod (10^9 + 7)$ (which would fit in an \texttt{int}) but
  computing the answer requires \emph{multiplying} $\bmod (10^9 + 7)$.
\end{warning}

\todo{Heap's Algorithm for generating all permutations.  See Bit
  Tricks for generating all subsets.}

\todo{PIE?}

\chapter{Bit Tricks} \label{chap:bittricks}

\todo{Basic bit manipulation.  Using bitstrings to compactly represent
  sets/states.  Iterating through all subsets with counter.}

\kattis{flipfive}

\todo{BitSet instead of array of booleans.}

\chapter{Geometry} \label{chap:geometry}

\todo{Points, vectors, angles.  Degrees/radians. \texttt{atan2}. Dot
  product. Rotation. Vector magnitude, norm (squared), normalize.
  Perpendicular (generate, test).}
\todo{Cross product in 2D. Signed area (parallelogram, triangle),
  polygon area, right/left turn, inside/outside testing.}
\todo{Lines/rays (point + vector).  Line intersection. Segment
  intersection. Closest point on a line/segment. Point/line distance.}
\todo{Convex hull.}

\chapter{Miscellaneous}

\section{2D grids}

2D grids/arrays (of characters, numbers, booleans\dots) are a popular
feature of many competitive programming problems.

\begin{itemize}
\item In many cases the grid should be thought of as a graph where
  each cell is a vertex which is connected by edges to its neighbors.
  Note that in these cases one rarely wants to explicitly construct
  a different representation of the graph, but simply use the grid
  itself as an (implicit) graph representation.
\item It is often useful to be able to assign a unique number to each
  cell in the grid, so we can store ID numbers of cells in data
  structures rather than making some class to represent a pair of a
  row and column index.  The easiest method is to number the first row
  from $0$ to $C-1$ (where $C$ is the number of columns), then the
  second row $C$ to $2C-1$, and so on.

  \begin{tabular}{|c|c|c|c|c|}
    \hline
    $0$ & $1$ & $2$ & $\dots$ & $C-1$ \\
    \hline
    $C$ & $C+1$ & $C+2$ & $\dots$ & $2C-1$ \\
    \hline
    $2C$ & $2C+1$ & $2C+2$ & $\dots$ & $3C-1$ \\
    \hline
    $\vdots$ & $\vdots$ & $\vdots$ & $\ddots$ & $\vdots$ \\
    \hline
    $(R-1)C$ & $(R-1)C+1$ & $(R-1)C+2$ & $\dots$ & $RC-1$ \\
    \hline
  \end{tabular}

\item Using this scheme, to convert between $(r,c)$ pairs and ID
  numbers $n$, one can use the formulas
  \[ (r,c) \mapsto r \cdot C + c \qquad n \mapsto (n / C, n \% C) \]

\item To list the four neighbors of a given cell $(r,c)$ to the north,
  east, south, and west, one can of course simply list the four cases
  manually, but sometimes this is tedious and error-prone, especially
  if there is a lot of code to handle each neighbor that needs to be
  copied four times.

  Instead, one can use the following template. The idea is that
  $(dr, dc)$ specifies the \emph{offset} from the current cell $(r,c)$
  to one of its neighbors; each time through the loop we rotate it
  counterclockwise by $1/4$ turn using the mapping
  $(dr,dc) \mapsto (-dc,dr)$ (see \fulllink{chap:geometry}).

  \code{grid/Neighbors.java}
\end{itemize}

\section{Range queries} \label{sec:range-queries}

Suppose we have a $1$-indexed array $A[1 \dots n]$ containing some
values, and there is some operation $\oplus$ which takes two values
and combines them to produce a new value.  Given indices $i$ and $j$,
we want to quickly find the value that results from combining all the
values in the range $A[i \dots j]$, \ie $A[i] \oplus A[i+1] \oplus
\dots \oplus A[j]$.

For example, $A$ could be an array of integers, and $\oplus$ could be
$\max$, that is, we want to find the maximum value in the range
$A[i \dots j]$.  Likewise $\oplus$ could be sum, or product, or GCD.
Or $A$ could be an array of booleans, and we want to find the AND, OR,
or XOR of the range $A[i \dots j]$.

\begin{itemize}
\item For this to make sense, the combining operation must typically
  be \emph{associative}, \ie $a \oplus (b \oplus c) = (a \oplus b) \oplus
  c$.  (This is called a \emph{semigroup}.)
\item Sometimes there is also an inverse operation $\ominus$ which
  ``cancels out'' the effects of the combining operation (this is
  called a \emph{group}).  For example, subtraction cancels out
  addition. On the other hand, there is no operation that can cancel
  out the effect of taking a maximum.
\item If we only need to find the value of combining a \emph{single}
  range $A[i \dots j]$, then ignore everything in this section and
  simply iterate through the interval, combining all the values in
  $O(n)$ time.
\item More typically, we need to do many queries, and $O(n)$ per query
  is not fast enough.  The idea is to preprocess the array into a data
  structure which allows us to answer queries more quickly, \ie in
  $O(1)$ or $O(\lg n)$.
\item Sometimes we also need to be able to \emph{update} the array in
  between queries; in this case we need a more sophisticated query
  data structure that can be quickly updated.
\end{itemize}

Each of the below subsections outlines one approach to solving this
problem; for quick reference, each subsection title says whether an
inverse operation is required, how fast queries are, and whether the
technique can handle updates.

\subsection{Prefix scan (inverse required; $O(1)$ queries; no updates)}

In a situation where we have an inverse operation and we do not need
to update the array, there is a very simple solution.  First, make
a \emph{prefix scan array} $P[0 \dots n]$ such that $P[i]$
stores the value that results from combining $A[1 \dots i]$.  ($P[0]$
stores the ``identity'' value, \eg zero if the combining operation is
sum.)  $P$ can be computed in linear time by scanning from left to
right; each $P[i] = P[i-1] \oplus A[i]$.  Now the value of $A[i \dots
j]$ can be computed in $O(1)$ time as $P[j] \ominus P[i-1]$.

Note that having $P[0]$ store the identity value is not strictly
necessary, but it removes the need for a special case.  If $A$ is
already $0$-indexed instead of $1$-indexed, then it's probably easier
to just put in a special case for looking up the value of $A[0 \dots
j]$ as $P[j]$, without the need for an inverse operation.

For example, suppose we are given an array of $10^5$ integers, along
with $10^5$ pairs $(i,j)$ for which we must output the sum of
$A[i \dots j]$.  Simply adding up the values in each range would be
too slow. We could solve this with the following code:

\code{range/PrefixSum.java}

More commonly, a prefix scan is a necessary first step in a more
complex solution. \kattis{divisible, dvoniz, srednji, subseqhard}

\subsection{Kadane's Algorithm}

As an aside, suppose we want to find the subsequence $A[i\dots j]$
with the \emph{biggest} sum.  A brute-force approach is $O(n^3)$:
iterate through all $(i,j)$ pairs and find the sum of each
subsequence.  Using the prefix scan approach, we can cut this down to
$O(n^2)$, since we can compute the sums of the $O(n^2)$ possible
subsequences in $O(1)$ time each.  However, there is an even better
$O(n)$ algorithm which is worth knowing, known as \emph{Kadane's
  Algorithm}.

The basic idea is simple: scan through the array, keeping a running
sum in an accumulator, and also keeping track of the biggest total
seen.  Whenever the running sum drops below zero, reset it to zero.
Below is a sample solution to \inlinekattis{commercials}.  Note that
subtracting \texttt{P} from each input is specific to the problem, but
the rest is purely Kadane's Algorithm.

\code{range/Commercials.java}

\subsection{2D prefix scan}

\todo{make pictures}

It is possible to extend the prefix scan idea to two dimensions.
Given a 2D array $A$, we create a parallel 2D array $P$ such that
$P[i][j]$ is the result of combining all the entries of $A$ in the
rectangle from the upper-left corner to $(i,j)$ inclusive.  The
simplest way to do this is to compute \[ P[i][j] = A[i][j] + P[i-1][j]
  + P[i][j-1] - P[i-1][j-1] \] Including $P[i-1][j]$ and $P[i][j-1]$
double counts all the entries in the rectangle from the upper left to
$(i-1,j-1)$ so we have to subtract them.

Given $P$, to compute the combination of the elements in some
rectangle from $(a,b)$ to $(c,d)$, we can compute \[ P[c][d] -
  P[a-1][d] - P[c][b-1] + P[a-1][b-1] \]

\inlinekattis{prozor} can be solved by brute force, but it's a nice
exercise to solve it using the above approach.

\subsection{Doubling windows (no inverse; $O(1)$ queries; no updates)}

\todo{Include link to discussion in CP3}

\subsection{Fenwick trees (inverse required; $O(\lg n)$ queries;
  $O(\lg n)$ updates)}

\kattis{fenwick, supercomputer, turbo, moviecollection}

We can use a \emph{Fenwick tree} to query the range $A[i..j]$ (\ie get
the combination of all the values in the range $A[i] \dots A[j]$
according to the combining operation $\oplus$) in $O(\lg n)$ time.  We
can also dynamically update any entry in the array in $O(\lg n)$ time.
If dynamic updates are required but we have an invertible combining
operation, a Fenwick tree should definitely be the first choice
because the code is quite short.

The code shown here stores \texttt{int} values and uses addition as
the combining operation, so range queries return the \emph{sum} of all
values in the range; but it can be easily modified for any other type
of values and any other invertible combining operation: change the
type of the array, change the \texttt{+} operation in the
\texttt{prefix} and \texttt{add} methods, change the subtraction
in the \texttt{range} method, and change the assignment \texttt{s = 0}
in \texttt{prefix} to the identity element instead of zero.

\begin{warning}
  Note that this \texttt{FenwickTree} code assumes the
  underlying array is $1$-indexed!
\end{warning}

\code{range/FenwickTree.java}

\begin{itemize}
\item The constructor creates a \texttt{FenwickTree} over an array of
  all zeros.
\item To create a \texttt{FenwickTree} over a given $1$-indexed array
  $A$, simply create a default tree and then loop through the array,
  calling \texttt{ft.add(i, A[i])} for each \texttt{i}.  This takes
  $O(n \lg n)$.
\item \texttt{ft.add(i, delta)} can be used to update the value at a
  particular index by adding \texttt{delta} to it.
\item If you want to simply replace the value at index $i$ instead of
  adding something to it, you could use \texttt{ft.add(i, newValue - ft.range(i,i))}.
\item \texttt{ft.range(i,j)} returns the sum $A[i] + \dots + A[j]$.
\end{itemize}

\todo{Discuss CP3 presentation of Fenwick trees; explain how Fenwick
  trees work}

\subsection{Segment trees (no inverse required; $O(\lg n)$ queries;
  $O(\lg n)$ updates)}

\todo{Segment trees.}

\chapter{Python} \label{chap:python}

Python's built-in support for arbitrary-size integers (using
\texttt{BigInteger} in Java is a pain!) and built-in dictionaries with
lightweight syntax make it attractive for certain kinds of problems.

Below is a basic template showing how to read typical contest problem
input in Python:

\pycode{python/template.py}

\todo{Mention basic Python data structures such as set, deque, list methods}

\end{document}
